\begin{ouverture}[algorithmique]{18}{Raisonner et préparer le brevet}{Partie 5 · Pensée informatique et raisonnement}
{Conjecturer · Démontrer · Réfuter · Rédiger · Vérifier}
\objectif{Distinguer exemple, conjecture, preuve et contre-exemple.}
\objectif{Utiliser réciproque, contraposée et raisonnement par l’absurde.}
\objectif{Organiser le temps et la rédaction de l’épreuve du brevet.}
\prerequis{Ensemble des chapitres précédents.}
\end{ouverture}
\pagestyle{otMain}\essentiel\fontsize{11}{15}\selectfont
\essSection{1}{Les formes du raisonnement}
\begin{definition}Une conjecture est une affirmation que l’on pense vraie après des observations. Une preuve établit une conclusion à partir d’hypothèses, de définitions et de propriétés. Un seul contre-exemple suffit à réfuter une affirmation générale ; plusieurs exemples ne suffisent pas à la démontrer.\end{definition}
\begin{definition}Pour une propriété « si $P$, alors $Q$ », la réciproque est « si $Q$, alors $P$ » et la contraposée « si $Q$ est faux, alors $P$ est faux ». La contraposée d’une propriété vraie est vraie ; sa réciproque peut être fausse.\end{definition}
\begin{savoirfaire}[Raisonner par l’absurde]Supposer le contraire de ce que l’on veut prouver, en déduire une contradiction, puis conclure que la supposition était fausse. Exemple : un triangle ne peut pas avoir deux angles droits, sinon la somme de ses angles dépasserait $180^\circ$.\end{savoirfaire}
\begin{savoirfaire}[Démontrer]Pour une propriété des entiers, traduire « pair » par $2n$, « impair » par $2n+1$, « multiple de $d$ » par $dn$. En géométrie, distinguer une définition, une propriété et une propriété caractéristique, puis citer les hypothèses avant le théorème.\end{savoirfaire}
\begin{attention}Une mesure sur le dessin, une image de calculatrice ou un essai de programme donne une observation ; cela ne justifie pas à soi seul une propriété pour tous les cas.\end{attention}
\clearpage
\essSection{2}{Une copie claire dans le temps imparti}
\begin{savoirfaire}[L’épreuve retenue]D’après la définition en vigueur à compter de la session 2027, retenue pour 2029, l’épreuve dure 2 h. Partie 1 : 20 min, sans calculatrice, 6 points d’automatismes et d’applications directes, à réponses courtes. Partie 2 : 1 h 40, calculatrice autorisée, 14 points de raisonnement et de résolution de problèmes. La première copie est relevée après 20 minutes.\end{savoirfaire}
\begin{propriete}La seconde partie ne comporte pas de QCM. Toutes les réponses doivent être justifiées, sauf indication contraire. La maîtrise de la langue est évaluée sur 2 points ; les essais et démarches engagés, même non aboutis, sont pris en compte.\end{propriete}
\begin{savoirfaire}[Rédiger]Écrire le calcul ou la propriété utile, donner les étapes essentielles, indiquer unités et arrondis, puis répondre par une phrase correctement construite. Si un résultat intermédiaire manque, poursuivre avec une valeur annoncée ou une démarche explicitée.\end{savoirfaire}
\begin{savoirfaire}[Vérifier]Relire les nombres copiés, le signe, la précision, les unités, les domaines et les ordres de grandeur. Une probabilité ne dépasse pas 1 ; une aire reste positive ; un achat entier se majore pour couvrir un besoin.\end{savoirfaire}
\clearpage\rubrique{Découverte}
\begin{blocLivre}{ex-18-1}\begin{exercice}[1]{Un contre-exemple}
« Le carré d’un nombre est toujours plus grand que ce nombre. » Réfuter cette phrase.
\end{exercice}\end{blocLivre}
\begin{blocLivre}{ex-18-2}\begin{exercice}[1]{Deux entiers impairs}
Démontrer que la somme de deux entiers impairs est paire.
\end{exercice}\end{blocLivre}
\begin{blocLivre}{ex-18-3}\begin{exercice}[1]{Réciproque et contraposée}
On considère la propriété vraie : « Si un entier est multiple de 6, alors il est pair. » Énoncer sa réciproque et sa contraposée. Lesquelles sont vraies ?
\end{exercice}\end{blocLivre}
\begin{blocLivre}{ex-18-4}\begin{exercice}[1]{Un arrondi utile}
Il faut 7,2 L d’un produit vendu en bidons de 3 L. Combien de bidons au minimum ?
\end{exercice}\end{blocLivre}
\begin{blocLivre}{ex-18-5}\begin{exercice}[1]{Une hypothèse manquante}
Un triangle est dessiné avec deux côtés de 3 cm et 4 cm. Peut-on affirmer que le troisième vaut 5 cm ?
\end{exercice}\end{blocLivre}

\clearpage\rubrique{Entraînement}
\begin{blocLivre}{ex-18-6}\begin{exercice}[2]{Un programme constant}
Choisir x, ajouter 5, doubler, puis retirer deux fois le nombre de départ. Montrer que le résultat est toujours le même.
\end{exercice}\end{blocLivre}
\begin{blocLivre}{ex-18-7}\begin{exercice}[2]{Un raisonnement par l’absurde}
Démontrer par l’absurde qu’un triangle rectangle ne peut pas avoir un angle obtus.
\end{exercice}\end{blocLivre}
\begin{blocLivre}{ex-18-8}\begin{exercice}[2]{Une conjecture de divisibilité}
On observe que $n(n+1)$ est pair pour plusieurs entiers n. Le démontrer.
\end{exercice}\end{blocLivre}
\begin{blocLivre}{ex-18-9}\begin{exercice}[2]{Vérifier une équation}
Un élève affirme que $x=4$ résout $3(x-2)=x+4$. Vérifier, puis résoudre.
\end{exercice}\end{blocLivre}
\begin{blocLivre}{ex-18-10}\begin{exercice}[2]{Une moyenne mal interprétée}
Une classe a une moyenne de 12/20. Est-il certain qu’au moins la moitié des notes sont supérieures ou égales à 12 ?
\end{exercice}\end{blocLivre}

\clearpage\rubrique{Problèmes}
\begin{blocLivre}{ex-18-11}\begin{exercice}[3]{Une commande au meilleur prix}
Pour 24 carnets : A vend à 2,80 € l’unité avec 6 € de livraison ; B vend des lots de 8 à 21 €, livraison gratuite. Comparer les coûts.
\end{exercice}\end{blocLivre}
\begin{blocLivre}{ex-18-12}\begin{exercice}[3]{Un modèle de bassin}
Un bassin pavé de 4 m sur 2,5 m contient 6000 L. La hauteur d’eau maximale est 0,8 m. Calculer la hauteur actuelle puis la quantité encore possible.
\label{fig-ex-18-12}\figureGeo{synthese-bassin}
\end{exercice}\end{blocLivre}
\begin{blocLivre}{ex-18-13}\begin{exercice}[3]{Un hasard et un budget}
Une tombola propose 200 billets à 3 €. Vingt billets gagnent chacun un lot coûtant 15 € à l’organisateur. Tous les billets sont vendus et chaque billet a la même chance. Calculer la probabilité qu’un billet donné gagne et le solde de l’organisation, sans autres frais.
\end{exercice}\end{blocLivre}
\begin{blocLivre}{ex-18-14}\begin{exercice}[3]{Un maximum expliqué}
Un rectangle a un périmètre de 40 cm. Montrer que son aire vaut $100-(x-10)^2$ si un côté vaut x, puis déterminer la plus grande aire.
\label{fig-ex-18-14}\figureGeo{rectangle-fonction}
\end{exercice}\end{blocLivre}
\begin{blocLivre}{ex-18-15}\begin{exercice}[3]{Un plan de révision}
Une élève réussit 8 automatismes sur 12, puis obtient 9/14 aux problèmes, rédaction comprise. Selon le barème des brevets blancs du livre, quelle note obtient-elle ? Quel travail cibler si ses quatre erreurs concernent les fractions ?
\end{exercice}\end{blocLivre}

\clearpage
